% discussion.tex -- Discussion and Conclusion (merged to save space)

The experiments validated the core claim: Langevin dynamics on the modern Hopfield energy converts deterministic attention into a principled stochastic sampler governed by a single temperature parameter.
The temperature spectrum (Fig.~\ref{fig:experiments}a) exhibited a smooth phase transition between retrieval and generation, the continuous analogue of the classical Hopfield/Boltzmann duality~\cite{ackeleyLearningAlgorithmBoltzmann1985} lifted to continuous states and Langevin dynamics. The convergence experiment (Fig.~\ref{fig:experiments}b) confirmed that Algorithm~1 reached the correct Boltzmann target, consistent with ULA theory~\cite{durmusNonasymptoticAnalysisUnadjusted2017}. The load-ratio sweep (Fig.~\ref{fig:phase-diagram}) showed that generation quality degraded gracefully as $K/d$ increased, providing empirical evidence linking modern Hopfield capacity to generation quality.

The strongest quantitative result came from protein sequence generation (Table~\ref{tab:protein-main}): on the Pfam RRM family ($K{=}68$, $d{=}59$), SA in the generation regime ($\beta{=}8$, $2\beta^*$) achieved $6.9{\times}$ lower global amino acid composition divergence than the VAE (KL${=}0.060$ vs.\ $0.416$) and $3.4{\times}$ lower per-position KL divergence ($2.92$ vs.\ $9.99$) at matched novelty ($0.623$ vs.\ $0.621$). SA-generated sequences also reproduce pairwise co-evolutionary coupling patterns with $r{=}0.871$ Pearson correlation to the stored family's MI matrix, compared to $r{=}0.525$ for the VAE. All generated sequences pass the Pfam RRM profile HMM at $E$-value~$<0.01$ (median $8.1{\times}10^{-34}$), confirming biological validity. These results demonstrate the core advantage of the training-free score function at multiple levels of biological structure: because the Langevin gradient is derived from the stored patterns at every step, it preserves family-level fidelity in position-specific conservation and pairwise coupling, not just global composition. The entropy inflection condition correctly predicted a lower $\beta^*{=}3.85$ (half the random-pattern value $\sqrt{d}{=}7.68$), reflecting the increased effective similarity variance introduced by conserved residues in the RRM motif.

On MNIST digit images (Table~\ref{tab:mnist-baselines}), SA at generation temperature led the VAE by $2.6{\times}$ on novelty and $2.0{\times}$ on diversity, while matching the Metropolis-corrected MALA baseline ($99.2\%$ acceptance rate at $\alpha{=}0.01$). A DDPM baseline failed entirely at $K{=}100$, producing samples indistinguishable from isotropic noise (max-$\cos{=}0.062$). A scaling study across $K\in\{100,\,500,\,1{,}000,\,2{,}000,\,3{,}500\}$ (Appendix~\ref{app:scaling}) confirmed that this failure persists across memory sizes: DDPM's max-cosine never exceeds $0.09$, while SA maintains $0.13$--$0.18$ in the generation regime and $0.24$--$0.33$ in the retrieval regime. Appendices~\ref{app:finance} and~\ref{app:simpsons} extended these results to S\&P~500 log-returns ($d{=}424$) and face images ($d{=}4{,}096$), confirming the SNR rule and the novelty--fidelity trade-off across a $5.2{\times}$ dimension range.

A deliberate choice shapes our experimental design: we evaluate at memory sizes $K$ ranging from 68 to 3{,}500, the low-to-moderate data regime where learned generative models are most constrained. This is not a limitation but the method's natural operating point. Score-based and diffusion models require enough training data to learn an accurate score function; when $K$ is small, this learning problem is ill-posed, and our DDPM baseline's failure across all $K$ tested (max-cosine $<0.09$) illustrates the consequence. Stochastic attention sidesteps this bottleneck entirely: the score function $\nabla\log p_\beta = \beta(\mathbf{T}-\boldsymbol{\xi})$ is exact for \emph{any} $K$, with fidelity that improves as the memory geometry becomes more structured. The protein experiment ($K{=}68$) is the clearest demonstration---the majority of Pfam families live in this regime, and no amount of architectural tuning will give a learned model access to training signal that does not exist. As $K$ grows, learned models will eventually close the gap; the scaling study (Appendix~\ref{app:scaling}) begins to map where this crossover occurs, and characterizing it precisely is an important direction for future work.

In summary, stochastic attention reuses the same primitive operations as a standard attention head and requires no learned score network, no training loop, and no contrastive objective. The analytic structure of the Hopfield energy (smoothness, dissipativity, and an exact closed-form gradient) yields a convergence guarantee in the convex regime ($\beta\sigma_{\max}^{2}<2$, Corollary~\ref{cor:convergence}) and geometric ergodicity of the continuous-time SDE for all $\beta>0$; at high $\beta$ the energy becomes multimodal and quantitative mixing bounds remain open, but the multi-chain protocol combined with fast intra-basin mixing provides a practical workaround (Section~4). The entropy inflection analysis (Proposition~\ref{prop:entropy-derivative}) provides a principled criterion for the retrieval--generation boundary: for random unit-norm memories the critical SNR scales as $d^{-1/4}$ (Eq.~\ref{eq:snr-star}), recovering the empirically observed threshold of $0.025$ at $d{=}64$. For structured data the effective variance of the similarities replaces $1/d$, but the inflection condition~\eqref{eq:inflection} remains a fully data-dependent selection rule for $\beta$.

\paragraph{Limitations.}
The ULA discretization introduces bias scaling with $\alpha$; our MALA comparison (Table~\ref{tab:mnist-baselines}) shows a 99.2\% acceptance rate at $\alpha{=}0.01$, confirming negligible bias at this step size. At larger $\alpha$, MALA (Algorithm~\ref{alg:mala}, Appendix~\ref{app:mala-algorithm}) is the preferred variant. At high temperature ($\beta{=}200$), SA samples are structured interpolations between stored patterns rather than crisp novel images: a Gaussian noise control (Appendix~\ref{app:noise-control}) confirms that SA retains meaningful structural similarity to stored patterns (max-$\cos=0.453$ vs.\ $0.328$ for variance-matched noise), but sample quality is limited by the dominance of the noise term over the gradient. Higher-fidelity generation requires operating closer to the phase-transition band. As discussed in Section~4, inter-basin mixing time at high $\beta$ grows exponentially with the barrier height; quantitative scaling laws as a function of $\beta$ and $K/d$ remain to be established. Each step costs $\mathcal{O}(dK)$ (two matrix--vector products plus a softmax), identical to a single attention head; for $d{=}784$, $K{=}100$ the full 30-chain experiment completes in under 10 minutes on a laptop CPU. The method also requires $\mathbf{X}$ to be known and fixed at sampling time; streaming or latent memory settings would require extension. A toy multi-head extension using PCA-partitioned subspaces (Appendix~\ref{app:multihead}) shows that the mechanism generalizes, but extension to learned multi-head projections, hierarchical memories, or full encoder-decoder transformers is future work.
